\subsection{End-to-End Effectiveness (RQ1)}
\label{subsec:e2e_evaluation}

We first evaluate whether SoCPilot can produce a functionally correct CPU--accelerator SoC and whether the resulting system improves complete-application latency over its paired CPU-only implementation. FPGA board execution is used as the final validation step. Importantly, this experiment evaluates the generated hardware and SoC configuration with manually prepared embedded runtime software; it does not attribute embedded-software generation to SoCPilot.

% Missing-result table: fill with measured board-level latency and validation outcomes.
\begin{table*}[t]
\centering
\caption{End-to-end effectiveness and FPGA-board validation of SoCPilot-generated SoCs.}
\label{tab:e2e_performance}
\renewcommand{\arraystretch}{1.12}
\resizebox{\textwidth}{!}{
\begin{tabular}{lllrrrcc}
\toprule
\textbf{Application} &
\textbf{Selected CPU} &
\textbf{Selected accelerator(s)} &
\textbf{CPU-only latency (ms)} &
\textbf{SoCPilot latency (ms)} &
\textbf{Speedup} &
\textbf{Correct} &
\textbf{Board validated} \\
\midrule
\texttt{plate\_recognition} & -- & -- & -- & -- & -- & -- & -- \\
\texttt{recognition}        & -- & -- & -- & -- & -- & -- & -- \\
\texttt{scan}               & -- & -- & -- & -- & -- & -- & -- \\
\texttt{wami}               & -- & -- & -- & -- & -- & -- & -- \\
\texttt{nightvision}        & -- & -- & -- & -- & -- & -- & -- \\
\midrule
Geometric mean                & -- & -- & -- & -- & -- & -- & -- \\
\bottomrule
\end{tabular}}
\end{table*}

Reported speedup is computed from complete-application latency rather than isolated accelerator latency. We therefore recommend reporting a latency breakdown that separates CPU-resident execution, accelerator execution, and CPU--accelerator communication/synchronization overhead for each benchmark.

% Recommended figure (data still required):
% Fig. X. End-to-end latency breakdown for CPU-only and SoCPilot designs.
% Components for SoCPilot: remaining CPU execution / accelerator execution /
% data transfer + invocation + synchronization.

Because the abstract reports the wall-clock completion time of the automated flow, the final evaluation should also document automation cost. The manually written embedded runtime is excluded from this timing and should be stated explicitly.

% Missing-result table: fill if the ``within 24 hours'' claim is retained.
\begin{table*}[t]
\centering
\caption{Automation cost of the SoCPilot hardware/SoC generation flow. Manually prepared embedded runtime software is excluded.}
\label{tab:automation_cost}
\renewcommand{\arraystretch}{1.12}
\begin{tabular}{lrrrr}
\toprule
\textbf{Application} &
\textbf{Total automated flow time (h)} &
\textbf{LLM calls} &
\textbf{Auto-repair iterations} &
\textbf{Manual edits within automated flow} \\
\midrule
\texttt{plate\_recognition} & -- & -- & -- & -- \\
\texttt{recognition}        & -- & -- & -- & -- \\
\texttt{scan}               & -- & -- & -- & -- \\
\texttt{wami}               & -- & -- & -- & -- \\
\texttt{nightvision}        & -- & -- & -- & -- \\
\bottomrule
\end{tabular}
\end{table*}

